\section{Introduction}
The extraction of hadron masses and matrix elements from
Monte Carlo estimates of Euclidean-space correlation functions
in lattice QCD can be done more reliably and accurately when
operators which couple more strongly to the states of interest
and less strongly to the higher-lying contaminating states
are used.  For states containing gluons, a crucial ingredient
in constructing such operators in lattice QCD is link
variable smearing.  Operators constructed out of
smeared or fuzzed links have dramatically reduced mixings with
the high frequency modes of the theory.  The use of such operators
has been shown to especially benefit determinations of the
glueball spectrum\cite{colin2}, hybrid meson masses\cite{hybrids,jkm1},
the torelon spectrum\cite{torelons}, and excitations of the static
quark-antiquark potential\cite{jkm2}.

Link variable smoothing is also playing an increasingly important
role in the construction of improved lattice actions.  In Ref.~\cite{fat1},
the use of so-called fat links in a staggered quark action was shown
to significantly decrease flavor symmetry breaking.
Smeared links were subsequently used\cite{fat2} to construct hypercubic
fermion actions having improved rotational invariance.
A staggered fermion action using a link smearing transformation\cite{hyp}
known as hypercubic (HYP) fat links was shown to improve flavor symmetry
by an order of magnitude relative to the standard action.
The so-called Asqtad improved staggered 
quark action\cite{asqtad1,asqtad2,asqtad3}
also makes use of link fattening to reduce flavor symmetry breaking.
Another variant of fermion actions which exploit smeared link variables
is the fat link irrelevant clover (FLIC) action\cite{flic}.
FLIC fermions are Wilson-like and described by an action which
includes an irrelevant clover improvement term constructed using
smeared links.  Fat links have also been used to construct a
gauge action\cite{fatgauge} with reduced discretization errors using
approximate renormalization group transformations.

The link fuzzing algorithm most often used in gluonic operator construction
is that described in
Ref.~\cite{APEsmear} in which every spatial link $U_j(x)$ on the lattice
is replaced by itself plus a real weight $\rho$ 
times the sum of its four neighboring
(spatial) staples, projected back into $SU(3)$.  Such a fuzzing step is
iterated $n_\rho$ times to obtain the final fuzzed link variables.
Empirically, one finds that the projection into $SU(3)$ is a crucial
ingredient of the smearing.  Link smearings which do not apply such
a projection are found to be much less effective.  The projection into
$SU(3)$ is not unique and must be carefully defined so that all symmetry
properties of the link variables are preserved.  Various ways of implementing
this projection have appeared in the literature.  Given a $3\times 3$
matrix $V$, its projection $U$ into $SU(3)$ is often taken to be the
matrix $U\in SU(3)$ which maximizes $\ReTr(UV^\dagger)$.  An iterative
procedure is required to perform such a maximization.  
Alternatively\cite{su3projectA}, the projected matrix $U$ can be defined by
$ U = V\ (V^\dagger V)^{-1/2} \det(V^{-1}V^\dagger)^{1/6}.$ 

Although the above smearing procedure works well in practice, the
projection is somewhat unpalatable since it may be viewed as an
arbitrary and abrupt way to remain within the group.  More importantly,
the branch cuts and lack of differentiability inherent in the projection
can hinder or even make impossible the application of modern Monte Carlo
updating techniques, such as hybrid Monte Carlo (HMC)\cite{hmc}, which
require knowing the response of the action to a small change in one of
the link variables.

A link smearing method which circumvents these problems is proposed
and tested in this paper.  The link smearing method is analytic
everywhere in the finite complex plane and utilizes the exponential
function in such a manner to remain within $SU(3)$, eliminating the
need for any projection back into the group.  Because of this construction,
the algorithm is useful for any Lie group.  The method is described
in Sec.~\ref{sec:analytic} and its practical implementations when
constructing operators and when computing the response of the action
to a change in a link variable are detailed in 
Secs.~\ref{sec:implement} and \ref{sec:force}, respectively.
Numerical tests of the smearing are presented in Sec.~\ref{sec:tests}.
Using the smeared plaquette and the effective mass $aE_{\rm eff}(t)$
associated with the static quark-antiquark potential, no
degradation in effectiveness is observed as compared to the standard
link smearing method, although an increased sensitivity to the
smearing parameter is found.  The results also suggest that lattice
actions and operators constructed out of smeared link variables may
be much less afflicted by radiative corrections since the
usually-dominant large tadpole contributions are drastically reduced.

